% !TEX program = platex
% 数学ノート集・共通テンプレート v2（積分論2026年度第1部の版面を参考）
% A5判。章内で定義・定理・例・注意等の番号を共通化。
% コンパイル： platex -interaction=nonstopmode math_note_template_v2.tex (2回)
%              dvipdfmx math_note_template_v2.dvi
\documentclass[a5paper,10pt,onecolumn]{jsarticle}
\usepackage[dvipdfmx,top=22mm,bottom=20mm,inner=19mm,outer=18mm,headheight=13pt,headsep=6mm]{geometry}
\usepackage{amsmath,amssymb,amsthm,mathtools}
\usepackage{enumitem,fancyhdr,titlesec}
\usepackage[dvipdfmx,bookmarks=true,colorlinks=false]{hyperref}
\usepackage{pxjahyper}

% ----- 差し替える文書情報 -----
\newcommand{\NoteTitle}{計算量理論の基礎}
\newcommand{\NoteSeries}{数学ノート集}

% ----- 文字・組版 -----
% 参考PDFの質素な明朝＋Computer Modernの印象を引き継ぎながら，
% A5判で窮屈にならないよう本文の行間を少し確保する．
\linespread{1.13}
\setlength{\parindent}{1zw}
\setlength{\parskip}{0pt}
\setlength{\abovedisplayskip}{6pt plus 2pt minus 2pt}
\setlength{\belowdisplayskip}{6pt plus 2pt minus 2pt}
\setlength{\abovedisplayshortskip}{3pt plus 2pt minus 1pt}
\setlength{\belowdisplayshortskip}{4pt plus 2pt minus 1pt}
\setlist[enumerate]{topsep=3pt,itemsep=1pt,parsep=0pt,leftmargin=2.4zw}
\titleformat{\section}{\normalfont\normalsize\bfseries}{\thesection}{0.8zw}{}
\titlespacing*{\section}{0pt}{2.0ex plus .4ex}{.9ex}
\titleformat{\subsection}{\normalfont\normalsize\bfseries}{\thesubsection}{0.7zw}{}
\titlespacing*{\subsection}{0pt}{1.5ex plus .3ex}{.6ex}
\raggedbottom

% ----- 柱とページ番号：ページ上部、罫線なし -----
\fancypagestyle{mathnote}{%
  \fancyhf{}%
  \fancyhead[L]{\footnotesize \NoteSeries\ （\NoteTitle）}%
  \fancyhead[R]{\footnotesize \thepage}%
  \renewcommand{\headrulewidth}{0pt}%
  \renewcommand{\footrulewidth}{0pt}%
}
\pagestyle{mathnote}

% ----- 節内通し番号の定義・定理群 -----
\newtheoremstyle{mathnotetheorem}
 {4.5pt}{4.5pt}{\normalfont}{}%
 {\bfseries}{}{0.65zw}%
 {\thmname{#1}\thmnumber{ #2}\thmnote{（#3）}}
\theoremstyle{mathnotetheorem}
\newtheorem{theorem}{定理}[section]
\newtheorem{definition}[theorem]{定義}
\newtheorem{lemma}[theorem]{補題}
\newtheorem{proposition}[theorem]{命題}
\newtheorem{corollary}[theorem]{系}
\newtheorem{example}[theorem]{例}
\newtheorem{remark}[theorem]{注意}
\renewcommand{\proofname}{証明}
\renewcommand{\qedsymbol}{$\square$}
\makeatletter
\renewenvironment{proof}[1][\proofname]{%
  \par\pushQED{\qed}\normalfont\topsep4\p@\@plus2\p@\relax
  \trivlist\item[\hskip\labelsep\bfseries #1]\ignorespaces
}{\popQED\endtrivlist\@endpefalse}
\makeatother

% ----- 使うノートで適宜拡張する記号 -----
\newcommand{\R}{\mathbb{R}}
\newcommand{\norm}[1]{\lVert #1\rVert}
\newcommand{\abs}[1]{\lvert #1\rvert}

\begin{document}
\section*{読了後にできること}
本稿を読むための前提は，計算量，最悪時計算時間，反転数．
本稿の具体的な到達課題は，具体的な整列算法の比較回数・交換回数を評価する．


\par\noindent\textbf{記号と約束.}\quad
集合の包含 $A\subseteq B$ は等号を許すものとする．
\par

\section{計算量の定義}
入力長 $n$ に対する決定的Turing機械の最悪計算時間を $T(n)$ とする．
\begin{definition}
$\mathbf P$ は決定的Turing機械で多項式時間内に決定できる言語の類，
$\mathbf{NP}$ は多項式 $p$ と多項式時間判定可能な関係 $R$ が存在して
\[
 x\in L\quad\Longleftrightarrow\quad
 \exists y\ (|y|\le p(|x|),\ R(x,y))
\]
と書ける言語 $L$ の類である．
\end{definition}
\begin{proposition}
$\mathbf P\subset\mathbf{NP}$ である．
\end{proposition}
\begin{proof}
$L\in\mathbf P$ の判定機械を $A$ とする．
証拠 $y$ を空文字列に固定し，$R(x,y)$ を $A(x)$ が受理することと定める．
検証は多項式時間である．
\end{proof}
\section{多項式時間帰着}
\begin{definition}
$A\le_p B$ とは，多項式時間で計算可能な全域写像 $f$ があり
$x\in A\iff f(x)\in B$ が成り立つことをいう．
\end{definition}
\begin{proposition}
$A\le_p B$，$B\in\mathbf P$ なら $A\in\mathbf P$．
また $A\le_p B$，$B\le_p C$ なら $A\le_p C$．
\end{proposition}
\begin{proof}
前半は $f$ の計算後に $B$ の判定器を実行すればよい．
後半は写像の合成を取る．多項式時間機械が出力する語の長さも
実行時間以下なので，多項式を合成した時間で計算できる．
\end{proof}
\begin{definition}
$B$ がNP困難であるとは，任意の $A\in\mathbf{NP}$ に対して
$A\le_p B$ が成り立つことをいう．
さらに $B\in\mathbf{NP}$ ならNP完全という．
\end{definition}
\begin{corollary}
NP完全な言語が $\mathbf P$ に属すれば $\mathbf P=\mathbf{NP}$ である．
\end{corollary}
\begin{proof}
任意の $A\in\mathbf{NP}$ は $A\le_p B$ を満たすので
前命題から $A\in\mathbf P$．逆包含は既に示した．
\end{proof}
\section{具体例}
\begin{proposition}
無向グラフのHamilton閉路の存在を判定する言語は $\mathbf{NP}$ に属する．
\end{proposition}
\begin{proof}
証拠として頂点の置換を与える．
各頂点がちょうど一回現れ，隣接する頂点と最後と最初の間に辺があるかを
多項式時間で確認できる．証拠長も頂点数の多項式である．
\end{proof}
\begin{remark}
NP完全性の証明には，任意のNP言語からの帰着が必要である．
特定の難しそうな入力例の存在だけでは足りない．
\end{remark}


\section{空間計算量とSavitchの定理}
\begin{definition}
$\mathrm{DSPACE}(s(n))$ と $\mathrm{NSPACE}(s(n))$ をそれぞれ
入力長 $n$ について $O(s(n))$ 個の作業テープセルを使用する
決定的・非決定的Turing機械が決定する言語の類とする．
入力テープは読み取り専用とし，機械はすべての計算枝で停止するとする．
\end{definition}
\begin{theorem}[Savitch]
$s(n)\ge\log_2 n$ が空間構成可能なら
\[
 \mathrm{NSPACE}(s(n))\subseteq\mathrm{DSPACE}(s(n)^2).
\]
\end{theorem}
\begin{proof}
非決定的機械の入力 $x$ を固定する．作業領域が $O(s)$ に限られるので，
状態・ヘッド位置・作業テープの内容からなる配置の個数は
$N\le 2^{Cs}$ と評価できる（入力ヘッド位置は $n$ 通りだが
$\log n\le s$ で吸収できる）．
$\mathrm{PATH}(C,D,i)$ を「配置 $C$ から $D$ へ
長さ高々 $2^i$ の遷移路がある」という命題とする．
$i=0$ では $C=D$ または一遷移で到達できるかを直接判定する．
$i\ge1$ では
\[
 \mathrm{PATH}(C,D,i)\iff
 \exists E\;\bigl(\mathrm{PATH}(C,E,i-1)
 \land\mathrm{PATH}(E,D,i-1)\bigr)
\]
を用いる．$E$ を全配置にわたって逐次列挙すれば決定的に判定できる．
再帰の深さは $\lceil\log_2 N\rceil=O(s)$ であり，
各階層で保存する配置 $C,D,E$ と列挙カウンタは $O(s)$ ビットで足りる．
従って再帰に必要な総作業空間は $O(s^2)$ である．
受理計算があれば同じ配置を二度訪れない受理路が存在するため，
長さは $N$ 以下に短縮できる．
初期配置からいずれかの受理配置への
$\mathrm{PATH}(C,D,\lceil\log_2N\rceil)$ を判定すれば受理判定ができる．
\end{proof}
\begin{corollary}
$\mathrm{PSPACE}=\mathrm{NPSPACE}$ である．
\end{corollary}
\begin{proof}
一方の包含は自明である．非決定的多項式空間機械には
ある $s(n)=n^c$ を用いて前定理を適用でき，$s(n)^2$ も多項式である．
\end{proof}

\section{決定的空間計算量の階層}
\begin{definition}
非減少の整数値関数 $s:\mathbb N\to\mathbb N$ が空間構成可能であるとは，
入力長 $n$ に対し $O(s(n))$ の作業領域内で $s(n)$ を二進表示で計算でき，
さらにその大きさに比例する領域を使用する操作を実行できることをいう．
ここでは $s(n)\ge n$ を仮定する．
\end{definition}
\begin{theorem}[決定的空間階層の初等形]
$s(n)\ge n$ を非減少で空間構成可能な関数とする．
ある言語 $L\in\mathrm{DSPACE}(s(n))$ が存在して，
任意の $f(n)=o(s(n))$ に対し
$L\notin\mathrm{DSPACE}(f(n))$ となる．
\end{theorem}
\begin{proof}
決定的Turing機械を有限二進語で符号化し，各機械には
任意に長いパディングを付けた符号が存在するようにする．
入力 $x$ の長さを $n$ とし，$x$ が機械 $M$ の有効な符号なら
次の対角機械 $D$ を実行する．
まず $s(n)$ を空間 $O(s(n))$ で構成し，$M(x)$ を万能機械で模倣する．
模倣で使う作業領域が $s(n)/2$ を超えるか，計算ステップ数が
$2^{s(n)}$ を超えたら停止して拒否する．
それより前に $M(x)$ が停止したときは受理・拒否を逆転して停止する．
無効な符号は直ちに拒否する．
ステップ数のカウンタは $s(n)+1$ ビットで足り，
符号の保持に必要な領域は高々 $n\le s(n)$ である．
従って $D$ は必ず停止し，作業領域は $O(s(n))$ に収まる．

$L=L(D)$ と置く．仮に $L$ を空間 $f(n)=o(s(n))$ で決定する
機械 $M$ が存在したとする．その固定した機械の符号を
パディングして任意に長い入力 $x$ にできる．
十分大きい $n=|x|$ に対して，$M(x)$ が使用する作業領域は
$s(n)/2$ より小さい．
状態数とテープ文字種類は固定されているから，
この空間での配置数は
$n\,2^{O(f(n))}$ 以下である．
$\log n+O(f(n))=o(s(n))$（$s(n)\ge n$）なので，
十分大きい $n$ では配置数は $2^{s(n)}$ 未満となる．
$M$ は決定機械だから同じ配置を繰り返さずに停止し，
その実行は $D$ の設定した時間・空間制限内で終わる．
従って $D(x)$ は $M(x)$ の出力を逆転させる．
これは $D$ と $M$ が同じ言語 $L$ を決定することに矛盾する．
\end{proof}
\begin{remark}
この定理は決定的空間の厳密な階層を示しているが，
時間階層定理で必要となる万能模倣の時間オーバーヘッドの評価を
含むものではない．
\end{remark}

\section{決定的時間階層定理の一つの具体形}
ここでは一テープ決定的Turing機械の時間計算量を用いる．
整数値関数 $t(n)\ge n$ が\emph{時間構成可能}であるとは，
入力 $1^n$ から $t(n)$ を計算し，時間 $O(t(n))$ で停止できることをいう．
\begin{lemma}[単一テープによる万能模倣]
ある一テープの万能機械 $U$ が存在し，
長さ $n$ の入力中に記された機械 $M$ の符号と実際の入力 $x$ に対し，
$M(x)$ の最初の $T\ge n$ ステップを
$O(T^2)$ 時間で模倣できる．
ここで定数は万能機械 $U$ のみに依存する．
\end{lemma}
\begin{proof}
万能機械のテープに，模倣する $M$ の固定長でない有限遷移表，
ヘッド位置・状態，および訪問済みのテープ区間を連続して記す．
$T$ ステップ以内に訪問し得る区間は長さ $O(T+n)=O(T)$ である．
遷移表の符号は入力長 $n$ 以下である．
各模倣ステップで遷移表と訪問済み区間をそれぞれ走査し，
現在の状態と読んだ文字から次の状態・書込み文字・移動方向を決める．
この作業は一テープ上でも $O(T)$ 時間で実行できるよう，
状態・ヘッド位置の印を区間内に保持し，必要に応じてテープ区間を延長する．
従って $T$ 回の模倣の総時間は $O(T^2)$ である．
\end{proof}
\begin{theorem}[時間階層の二次分離形]
$t(n)\ge n$ を時間構成可能で非減少な関数とする．
ある言語 $L$ が存在して
\[
 L\in\operatorname{DTIME}(t(n)^2),\qquad
 L\notin\operatorname{DTIME}(o(t(n)))
\]
を満たす．時間構成可能性を満たす $t$ の選択により，
多項式時間クラスの内部にも厳密な階層が存在する．
\end{theorem}
\begin{proof}
決定的機械の符号は，意味を変えずに任意に長くパディングできるとする．
長さ $n$ の入力 $x$ が機械 $M$ の有効な符号なら，
対角機械 $D$ は $M(x)$ を高々 $t(n)$ ステップだけ模倣する．
制限以内に $M$ が受理すれば $D$ は拒否し，拒否すれば受理する．
停止しなければ $D$ は拒否する．無効な符号も拒否する．
時間構成可能性と前補題から $D$ の計算時間は $O(t(n)^2)$ である．

$L=L(D)$ と置く．仮に同じ $L$ を時間 $f(n)=o(t(n))$
で決定する機械 $M$ が存在したとする．
$M$ の符号を長くパディングした入力 $x$ を選べば，
十分大きい $n=|x|$ に対し $M(x)$ は $t(n)$ ステップ以内に停止する．
ところが $D(x)$ は $M(x)$ の受理・拒否を必ず逆転する．
$M$ と $D$ が同じ言語を決定することに矛盾する．
従って所要の分離が成立する．
\end{proof}
\begin{remark}
これは明示的な万能模倣の二次オーバーヘッドを用いた時間階層定理である．
標準的なより鋭い $t(n)\log t(n)$ 程度の分離には，
一テープから多テープへの模倣も含む高精度のシミュレーションが必要であり，
ここでは主張しない．
\end{remark}


\section{到達課題の完全解答}
\noindent\textbf{挿入ソートの最悪時計算量．}
\begin{proposition}
隣接交換を使う挿入ソートにおいて，
交換回数は入力順列の反転数に等しい．
逆順に並んだ長さ $n$ の相異なる数列では
交換回数は $n(n-1)/2$ となる．
\end{proposition}
\begin{proof}
一回の隣接交換は，逆順に並ぶ隣接二元を
正しい順序に入れ替えるので，
その二元の反転だけを解消し，
反転数をちょうど1減らす．
整列完了後の反転数は零だから，
最初の反転数と交換回数は等しい．
逆順の順列では全ての
$\binom n2$ 組が反転を成すので
交換回数は $\binom n2=n(n-1)/2$ である．
\end{proof}
\noindent\textbf{到達確認．}
入力 $(3,1,2)$ の反転数を求めよ．
\par\noindent\textbf{解答．}
$(3,1),(3,2)$ の二組なので2である．
隣接交換による挿入ソートは二回の交換で整列する．
\end{document}
